\addcontentsline{toc}{subsubsection}{Solution 10. The Reciprocal Pinch}

\begin{dlsummary}[title={\bfseries Solution 10. The Reciprocal Pinch},breakable]

\begin{proof}
The reciprocal identity gives $f(1)=1$. By superadditivity,
\[
f(2)\geq 2f(1)=2.
\]
Writing $a=f(2)$, we have $a\geq2$ and $f(1/2)=1/a$. Thus
\[
a+\frac1a\leq\frac52,
\]
or $(2a-1)(a-2)\leq0$. Since $a\geq2$, this forces
\[
f(2)=2. \tag{1}
\]

Applying the multiplicative inequality to $1/x,1/y$ and using the
reciprocal identity gives
\[
\frac1{f(xy)}=f\left(\frac1{xy}\right)
\geq f\left(\frac1x\right)f\left(\frac1y\right)
=\frac1{f(x)f(y)}.
\]
Together with the original inequality, this proves
\[
f(xy)=f(x)f(y). \tag{2}
\]
Also, $f$ is strictly increasing: if $x<y$, then
\[
f(y)=f(x+(y-x))\geq f(x)+f(y-x)>f(x).
\]

By (2) and (1), $f(2^n)=2^n$ for every integer $n\geq0$. Fix $t>0$,
and for sufficiently large $n$ put
\[
m_n=\left\lfloor\frac{2^n}{t}\right\rfloor.
\]
Since $m_nt\leq2^n$, monotonicity and repeated superadditivity give
\[
2^n=f(2^n)\geq f(m_nt)\geq m_nf(t).
\]
As $2^n/m_n\to t$, we obtain $f(t)\leq t$. Applying this bound to
$1/t$ and using the reciprocal identity yields
\[
f(t)=\frac1{f(1/t)}\geq t.
\]
Hence $f(t)=t$ for every $t>0$.
\end{proof}

The identity function directly satisfies all the conditions, so it is
the unique solution.

\end{dlsummary}